% Ad Atticum 12.43 --- headnote
% ad-atticum-12-43, 12 May 45 BC, Astura

\textit{Cicero to Atticus, written from Astura on the
fourth day before the Ides of May 709 AUC --- 12 May
45 BC (the manuscript dateline: \textit{Scr.\ Asturae
iv Id.\ Mai.\ a.\ 709 (45)}). A brief two-section
note, written as Cicero prepares the move out of
Astura. He approves of Atticus' decision to handle
his business from home --- evidently a domestic
withdrawal to take the affair through without the
interruptions of the city --- and confirms his own
schedule: the day after the Ides at Lanuvium, then
either to Rome or to the Tusculanum.}

\textit{The middle of the first section carries one
of the more disturbed manuscript passages in the
Astura sequence; the daggered text is preserved as
the editors mark it, and the sense returns at
\textit{believe me, to a degree you cannot imagine}.
What is clear is Cicero's admission that he knows
Atticus does not approve of the project (the shrine,
the gardens, the whole obsession), but that he is
bound to it past the point at which Atticus can be
expected merely to tolerate it: \textit{ferendus?
immo vero etiam adiuvandus}, ``bear with? no, indeed
even help.'' The second section runs through the
estate alternatives: Otho is failing, Clodia is the
next best, and Atticus is asked to bring off anything
he can --- including the Trebonian heirs --- before
the summer slips away. The vow Cicero claims to be
bound by is the inward one, the promise of the shrine
to Tullia.}
